% Summer research report: two-particle Hamiltonian system
% Data sources: report01-04.ipynb, case3_newton_failure_diagnosis.ipynb,
% lyapunov_*.py; orbit figures from report/make_orbit_figures.py.
% Text marked with \todo{...} still needs checking or a decision.
% Numbers were copied from notebook outputs on 2026-11-07; the equilibrium
% energy/frequencies and the step-halving check were computed the same day
% (see the comments next to them).

\documentclass[11pt,a4paper]{article}

\usepackage[margin=2.5cm]{geometry}
\usepackage{amsmath,amssymb,amsthm}
\usepackage{graphicx}
\usepackage{booktabs}
\usepackage{xcolor}
\usepackage[colorlinks=true,linkcolor=blue!50!black,citecolor=blue!50!black,urlcolor=blue!50!black]{hyperref}

\graphicspath{{../images/}{figures/}}

% Figures and tables use [H] (float package): placed exactly where they are
% in the source, right after the paragraph that first refers to them.
\usepackage{float}

% Visible placeholder. For a clean draft use \newcommand{\todo}[1]{}
\newcommand{\todo}[1]{\textcolor{red}{[TODO: #1]}}

\newcommand{\q}{\mathbf{q}}
\newcommand{\p}{\mathbf{p}}
\newcommand{\M}{\mathcal{M}}

\title{Periodic Orbits and Chaos in a System of Two Repelling Particles in a Box:\\
a Numerical Study}
\author{Jiajun Wu\\[2pt]
\small Department of Mathematics, Imperial College London\\
\small Summer research project (2026), supervised by Prof.\ Dmitry Turaev}
\date{\todo{date}}

\begin{document}
\maketitle

% ---------------------------------------------------------------------------
\begin{abstract}
We study numerically a Hamiltonian system of two equal particles in a square
box, repelling each other through an inverse-power potential and confined by
smooth, singular wall potentials. Combining a symplectic integrator, Poincar\'e
sections, Newton's method for fixed points of the Poincar\'e map, and finite-time
Lyapunov exponents, we find that the motion is regular close to the equilibrium,
becomes chaotic at slightly higher energies, and is again largely organised by
stable periodic motion at high energy. We compute two families of periodic
orbits and follow them in energy. A ``vertical'' bouncing orbit is linearly
stable for $E\gtrsim3.9$ and loses stability between $E\approx3.9$ and $3.3$,
where the numerics are consistent with a symmetry-breaking bifurcation. A
``circular'' orbit around the centre of the box is linearly stable at every
energy we computed, from $E\approx2$ to $E\approx202$. At $E\approx52$ the phase
space is mixed, with a chaotic sea and many elliptic islands whose centres we
locate. We also describe where our continuation method fails and why.
\end{abstract}

% ---------------------------------------------------------------------------
\section{Introduction}\label{sec:intro}

Consider two particles in a box that repel each other. Nothing could be simpler
to write down, yet the long-time behaviour of such a system is far from obvious.
The particles might settle into a regular, repeating pattern, or wander
chaotically over all the available space, or do one or the other depending on
the energy and the initial condition.

For a single particle in a smooth convex box, this is the classical question of
billiard-like dynamics. With two or more particles the system has more degrees
of freedom and, in general, no conserved quantity besides the energy, so one
expects coexistence of regular and chaotic motion. Rom-Kedar and Turaev
\cite{turaev2022} proved that $N$ particles with a smooth repelling interaction
in a compact region are \emph{not} ergodic at sufficiently high energy: there
are stable ``choreographic'' motions, in which the particles follow
approximately the same path close to an elliptic periodic orbit of the
single-particle problem. \todo{Check this summary and the framing of the next
sentence with Prof.\ Turaev.} This project explores numerically the simplest
case, two particles in a square, and asks what happens at energies that are
not necessarily large.

Concretely, we
(i)~integrate typical trajectories (Section~\ref{sec:trajectories}),
(ii)~use Poincar\'e sections to separate regular from chaotic motion as the
energy changes (Section~\ref{sec:sections}),
(iii)~compute periodic orbits with Newton's method, determine their linear
stability and follow them in energy
(Sections~\ref{sec:vertical}--\ref{sec:circular}), and
(iv)~map the phase space at a fixed high energy with finite-time Lyapunov
exponents (Section~\ref{sec:lyapunov}).
The model and its symmetries are described in Section~\ref{sec:model} and the
numerical methods in Section~\ref{sec:methods}. Section~\ref{sec:discussion}
discusses limitations and open questions.
All statements below are based on numerical evidence; none of them is proved.

% ---------------------------------------------------------------------------
\section{The model}\label{sec:model}

Two particles with positions $\q_i=(x_i,y_i)$ and momenta $\p_i$ ($i=1,2$)
move in the box $(0,L_x)\times(0,L_y)$ with Hamiltonian
\begin{equation}\label{eq:H}
  H(\q,\p)=\sum_{i=1}^{2}\frac{|\p_i|^2}{2m_i}
  + V\bigl(|\q_1-\q_2|\bigr) + U(\q_1) + U(\q_2),
\end{equation}
where the particle--particle repulsion and the wall potential are
\begin{equation}\label{eq:potentials}
  V(r)=k\,r^{-\alpha},\qquad
  U(x,y)=C\Bigl[x^{-\beta}+(L_x-x)^{-\beta}+y^{-\beta}+(L_y-y)^{-\beta}\Bigr].
\end{equation}
Both terms are singular: $V\to\infty$ at collision and $U\to\infty$ at the
walls. Hence, on any energy level, the particles stay a positive distance from
each other and from the walls, and the motion is defined for all time.

Throughout we take $k=C=\alpha=\beta=1$, $m_1=m_2=1$ and $L_x=L_y=10$.
The model is a dimensionless toy model, so these values were chosen to be of
order one rather than to match a physical system.
\todo{Mention whether other exponents or rectangular boxes were tried
(e.g.\ \texttt{test2.py}).}

\paragraph{Symmetries.}
The energy is conserved. Because of the walls, neither total momentum nor
angular momentum is conserved, and we are not aware of any further integral.
The system is invariant under the symmetries of the square applied to both
particles, and under exchange of the particles. Combining exchange with the
point reflection $\q\mapsto2\mathbf{c}-\q$ through the centre
$\mathbf{c}=(L_x/2,L_y/2)$ shows that the \emph{centre-symmetric set}
\begin{equation}\label{eq:centresym}
  \M=\bigl\{\q_2 = 2\mathbf{c}-\q_1,\ \p_2=-\p_1\bigr\}
\end{equation}
is invariant under the flow: it is the fixed-point set of the symmetry
$(\q_1,\q_2,\p_1,\p_2)\mapsto(2\mathbf{c}-\q_2,\,2\mathbf{c}-\q_1,\,-\p_2,\,-\p_1)$,
which preserves $H$. On $\M$ the motion is described by particle~1 alone, as a
system with two degrees of freedom.

\paragraph{Dimensions.}
The full phase space is $8$-dimensional, an energy level is $7$-dimensional,
and a Poincar\'e section inside an energy level is $6$-dimensional. On $\M$
these numbers become $4$, $3$ and $2$, so the Poincar\'e map restricted to $\M$
is an area-preserving map of a planar region, like the standard map. This is
what makes the two-dimensional pictures in Section~\ref{sec:lyapunov}
possible. All initial conditions in this report lie on $\M$.

\paragraph{Equilibrium.}
On the diagonal, $\nabla(V+U)=0$ has the symmetric solution
$\q_1=(a,a)$, $\q_2=(10-a,10-a)$ with $a\approx3.6645$, at energy
$E_{\mathrm{eq}}\approx1.9877$.
% Computed 2026-11-07 from the Hessian of V+U (SymPy), not in the notebooks.
Its linearisation has frequencies $\omega\approx0.107$, $0.220$ (double) and
$0.350$.
% Checked twice on 2026-11-07: SymPy Hessian and an independent
% finite-difference Hessian agree to 8 digits.
The modes with $\omega\approx0.107$ (motion perpendicular to the
diagonal) and $\omega\approx0.350$ (motion along the diagonal) lie in $\M$;
the double frequency belongs to the two centre-of-mass modes,
$\delta\q_1=\delta\q_2$, which leave $\M$.

% ---------------------------------------------------------------------------
\section{Numerical methods}\label{sec:methods}

\subsection{Time integration}
The equations of motion are integrated with the leapfrog (St\"ormer--Verlet)
scheme
\begin{equation}\label{eq:leapfrog}
  \p^{n+1/2}=\p^n+\tfrac h2 F(\q^n),\quad
  \q^{n+1}=\q^n+h\,M^{-1}\p^{n+1/2},\quad
  \p^{n+1}=\p^{n+1/2}+\tfrac h2 F(\q^{n+1}),
\end{equation}
where $F=-\nabla(V+U)$ is obtained by symbolic differentiation (SymPy) and
compiled into a NumPy function. Leapfrog is second order, symplectic and
time-reversible, so the energy error oscillates but does not drift over long
times \cite{hairer2006}. This matters here because we integrate for hundreds
of time units and want to trust what we see in Poincar\'e sections.

The step size is adapted to the speed of the particles: $h=10^{-2}$ near the
equilibrium, $h=2.5\times10^{-3}$ to $5\times10^{-4}$ for the periodic-orbit
computations, and down to $h=10^{-5}$ at $E\approx202$. In all runs used for
Poincar\'e sections the relative energy error at the crossings is at most
about $10^{-6}$. As a convergence check, halving $h$ twice changes the computed
periods of the orbits in Figures~\ref{fig:vertical} and~\ref{fig:circular} by
less than $10^{-6}$ in relative terms, with the factor-four reduction expected
of a second-order method, and the orbits still close to within $10^{-7}$.
% Step-halving check run 2026-11-07: E=11.075 T=6.87461211/...275/...291;
% E=3.876 T=14.80128691/...697/...699; E=3.011 T=23.80233268/...379/...407.

\subsection{Poincar\'e sections}
A Poincar\'e section is a hypersurface $\Sigma=\{g(\q)=0\}$, crossed in a fixed
direction; we use $y_1=5$ with $p_{y_1}>0$ and the diagonal $x_1=y_1$ with
$p_{x_1}-p_{y_1}>0$. A leapfrog step will generally jump over $\Sigma$ rather
than land on it, so the crossing point has to be reconstructed. We use two
methods, depending on what the crossing points are needed for.

\paragraph{Linear interpolation (for plotting sections).}
If $s_n=g(\q^n)<0\le s_{n+1}=g(\q^{n+1})$, we set
$\theta=-s_n/(s_{n+1}-s_n)$ and take
$(\q,\p)=(\q^n,\p^n)+\theta\bigl((\q^{n+1},\p^{n+1})-(\q^n,\p^n)\bigr)$ at time
$(n+\theta)h$. The error is $O(h^2)$, the same order as that of leapfrog
itself, which is more than enough for the section plots of
Figures~\ref{fig:cases}, \ref{fig:energy-sections} and~\ref{fig:lyapunov}.

\paragraph{Partial leapfrog steps (for Newton's method).}
To solve $P^k(z)=z$ we need a return map that is accurate and smooth in $z$,
because its Jacobian is computed by finite differences. Here the crossing is
found by taking a leapfrog step of length $\tau\in(0,h]$ from $(\q^n,\p^n)$
and solving $g(\q(\tau))=0$ for $\tau$: by bisection in the continuation runs
of Sections~\ref{sec:vertical}--\ref{sec:circular}, and by a few secant
iterations in the island-centre search of Section~\ref{sec:lyapunov}. The
crossing point then lies exactly on a (shortened) leapfrog trajectory and
satisfies $g=0$ to round-off, instead of only approximately.

\subsection{Periodic orbits by Newton's method}\label{sec:newton}
At fixed energy $E$, a point of $\Sigma$ is described by six coordinates $z$;
the remaining momentum component is recovered from $H=E$, so every iterate
stays on the energy level. A periodic orbit that crosses $\Sigma$ $k$ times per
period is a fixed point of $P^k$, where $P$ is the first-return map. We solve
$P^k(z)=z$ by Newton's method,
\begin{equation}\label{eq:newton}
  \bigl(DP^k(z_j)-I\bigr)\,\Delta z_j=-\bigl(P^k(z_j)-z_j\bigr),\qquad
  z_{j+1}=z_j+s_j\,\Delta z_j,
\end{equation}
where $DP^k$ is approximated by central finite differences and a backtracking
line search chooses $s_j\in(0,1]$ so that the residual decreases. From a
reasonable initial guess Newton typically converges in 3--6 iterations to a
residual $|P^k(z)-z|$ between $10^{-9}$ and $10^{-13}$.

\paragraph{Linear stability.}
The eigenvalues $\lambda$ of $DP^k(z^*)$ are the non-trivial Floquet
multipliers of the periodic orbit. Because the map is symplectic, they come in
quadruples $\lambda,\bar\lambda,1/\lambda,1/\bar\lambda$. We call the orbit
\emph{linearly stable} (elliptic) if all multipliers lie on the unit circle,
and \emph{unstable} if some $|\lambda|>1$. Linear stability does not imply
stability of the nonlinear system, although for an elliptic orbit KAM theory
makes nonlinear stability typical in two degrees of freedom
\cite{meiss}. \todo{Keep this sentence only if Prof.\ Turaev agrees with it.}
With finite-difference Jacobians, $|\lambda|$ is accurate to about $10^{-6}$,
so we treat $\bigl||\lambda|-1\bigr|<10^{-3}$ as being on the unit circle.

\todo{If keeping the transverse-stability calculation of
\texttt{lyapunov\_section\_periodic\_orbits.py}, explain it here:
perturbations leaving $\M$ with $\delta\q_2=\delta\q_1$ reduce to two Hill
equations, and $|\operatorname{tr}|<2$ for both gives stability in the full
phase space. The results are not stored, so the script would need to be
re-run.}

\subsection{Continuation in energy}
To follow a periodic orbit as the energy changes, we take a converged orbit at
energy $E$, rescale its momenta to the new energy $E'$, and use the result as
the initial guess for Newton at $E'$. This is natural-parameter continuation:
simple, but unable to pass a fold, and liable to jump to a nearby orbit when two
branches come close (Section~\ref{sec:discussion}).

\subsection{Finite-time Lyapunov exponents}
On $\M$ we integrate, together with each trajectory, a tangent vector updated
by the exact linearisation of the leapfrog step, and renormalise it every 40
steps. The finite-time maximal Lyapunov exponent $\lambda_{\max}$ is the
average logarithmic growth rate over $T=500$ ($h=10^{-3}$). Since all crossings
of one trajectory share its $\lambda_{\max}$, one integration colours all its
points in the Poincar\'e section. Initial conditions form a $120\times120$ grid
on the section, of which the 13\,920 energetically accessible points are used.
To locate island centres we search for fixed points of $P^k$, $k\le4$, starting
Newton from grid points that nearly return after $k$ crossings, and classify a
fixed point as elliptic if $|\operatorname{tr}DP^k|<2$.

% ---------------------------------------------------------------------------
\section{Results}\label{sec:results}

\subsection{Typical trajectories}\label{sec:trajectories}
Table~\ref{tab:cases} lists five initial conditions, all on $\M$, and
Figure~\ref{fig:cases} shows the trajectories together with 200 returns to a
Poincar\'e section. Cases~1 and~2 start near the equilibrium and stay near it:
the section points lie on smooth curves, so the motion is regular
(quasi-periodic). In Case~3 the particles move around the centre of the box,
always opposite each other, and the section points form two tiny clusters,
which suggests a periodic orbit nearby; it is computed in
Section~\ref{sec:circular}. In Case~4 each particle bounces between the top and
bottom walls along an almost vertical line. It is a periodic orbit: after one
period $T\approx1.0027$ the state returns to its initial value to within
$6\times10^{-7}$. Case~5 starts from the same momenta at a different $x$, and
its section points lie on a closed curve around the Case~4 point, as expected
near an elliptic periodic orbit.

\begin{table}[H]
\centering
\caption{Initial conditions. Particle~2 always starts at the point reflection
$2\mathbf{c}-\q_1(0)$ with momentum $-\p_1(0)$; $a\approx3.6645$ is the
equilibrium position, at energy $E_{\mathrm{eq}}\approx1.9877$.}
\label{tab:cases}
\begin{tabular}{@{}lllll@{}}
\toprule
Case & $\q_1(0)$ & $\p_1(0)$ & $E$ & Motion\\
\midrule
1 & $(a,a)+5\cdot10^{-4}(1,1)$ & $(10^{-3},0)$ & $1.9877$ & small regular oscillation\\
2 & $(a,a)+0.25\,(1,1)$ & $(0,0.05)$ & $2.0073$ & regular, quasi-periodic\\
3 & $(a,a)+5\cdot10^{-4}(1,1)$ & $(0.1,-0.1)$ & $2.0077$ & near a periodic loop\\
4 & $(3.8352,\,5)$ & $(0,20)$ & $402.08$ & periodic, vertical bouncing\\
5 & $(3,\,5)$ & $(0,20)$ & $402.00$ & regular, around Case 4\\
\bottomrule
\end{tabular}
\end{table}

\begin{figure}[H]
  \centering
  \includegraphics[width=\textwidth]{fig_cases.pdf}
  \caption{Top: trajectories for the five cases of Table~\ref{tab:cases}
  (particle~1 blue, particle~2 orange). Bottom: 200 returns to a Poincar\'e
  section, plotted relative to the starting point (red star). The section is
  the hyperplane through the starting point normal to the largest initial
  velocity component, crossed in the same direction; the axes are the other
  coordinate of the same particle and its momentum. Note the scales: in Case~4
  all returns lie within $10^{-5}$ of the start.}
  \label{fig:cases}
\end{figure}

\subsection{Regular and chaotic motion as the energy increases}\label{sec:sections}
To see how the type of motion depends on energy, we start from the Case~2
position and give the particles vertical momenta $\pm v$ for four values of $v$,
so that $E$ ranges from just above $E_{\mathrm{eq}}$ to $E\approx11$
(Figure~\ref{fig:energy-sections}). At $E\approx2.007$ the section points lie on
smooth curves and the motion is regular. Already at $E\approx2.028$ the points
scatter over two regions, and at $E\approx2.97$ the trajectory fills most of
the box while the section points are spread everywhere: the motion appears
chaotic. At $E\approx11$ the motion is regular again: the particles stay in two
thin vertical bands and the section points lie on a single closed curve,
because the orbit is now trapped near the stable vertical periodic orbit of
Section~\ref{sec:vertical}. So in this sequence chaos is strongest at
intermediate energies, and at high energy regular motion reappears around
stable periodic orbits, in line with \cite{turaev2022}.

\begin{figure}[H]
  \centering
  \includegraphics[width=\textwidth]{fig_energy_sections.pdf}
  \caption{Starting from the Case~2 position with vertical momenta
  $\pm v$: trajectories over $300$ time units (top) and 200 returns to the
  section $y_1=y_{\mathrm{eq}}$, $p_{y_1}>0$, in the coordinates
  $(x_1,p_{x_1})$ (bottom).}
  \label{fig:energy-sections}
\end{figure}

\subsection{The vertical bouncing family}\label{sec:vertical}
Starting from the Case~4 orbit with smaller momentum ($E\approx11.08$), we
follow the vertical orbit down in energy using the section $y_1=5$
(Figure~\ref{fig:vertical}). Newton converges at every step. For
$E\approx11.1$ down to $E\approx3.88$ all six multipliers lie on the unit
circle, so the orbit is linearly stable. At $E\approx3.32$ it is unstable:
one pair of multipliers has become real, $\lambda\approx2.59$ and $0.386$, and
the instability grows as $E$ decreases further ($\max|\lambda|\approx6.04$ at
$E\approx2.77$ and $5.21$ at $E\approx2.66$). The pair that leaves the unit circle
is at $0.70\pm0.72i$ at $E\approx3.88$, so it must pass through $+1$ between
$E\approx3.88$ and $3.32$. Figure~\ref{fig:multipliers}(a) shows this in
terms of the stability index $\rho=(\lambda+\lambda^{-1})/2$ of each pair,
which satisfies $|\rho|\le1$ exactly when the pair is on the unit circle. One
index also comes within $10^{-2}$ of $-1$ near $E\approx6.1$ and $E\approx3.9$.
Since our energy steps are about $0.55$, we cannot exclude short intervals of
instability (period-doubling) there.

% Generated by report/make_orbit_figures.py
\begin{figure}[H]
  \centering
  \includegraphics[width=\textwidth]{fig_vertical_family.pdf}
  \caption{One period of the vertical family. (a)--(b) linearly stable, with
  periods $T\approx6.87$ and $14.80$; (c) unstable, $\max|\lambda|\approx2.59$;
  (d) at $E\approx2.55$ Newton converges to a different, S-shaped orbit, which
  is linearly stable. The black dot is the section point ($y_1=5$,
  $p_{y_1}>0$).}
  \label{fig:vertical}
\end{figure}

% Generated by report/make_multiplier_figure.py
\begin{figure}[H]
  \centering
  \includegraphics[width=\textwidth]{fig_multipliers.pdf}
  \caption{Stability indices $\rho=(\lambda+\lambda^{-1})/2$ of the three
  multiplier pairs against energy; the shaded band $|\rho|\le1$ is linear
  stability. (a) Vertical family: one index leaves the band through $+1$
  between $E\approx3.88$ and $3.32$. Open squares: the S-shaped orbit found
  at $E\approx2.55$. (b) Circular family: all indices stay in the band, but one
  comes very close to $+1$ at $E\approx2.11$ and to $-1$ at $E\approx2.31$.
  In (b) two of the three pairs coincide.}
  \label{fig:multipliers}
\end{figure}

A multiplier pair colliding at $+1$ is where a symmetric periodic orbit can
lose stability through a pitchfork (symmetry-breaking) bifurcation, in which
two non-symmetric periodic orbits branch off. Our data are consistent with this:
at $E\approx2.55$, Newton no longer converges to the vertical orbit but to a
different, linearly stable orbit with $p_x\neq0$, whose paths are S-shaped
(Figure~\ref{fig:vertical}(d)). We have not resolved the bifurcation itself,
which would need finer energy steps near $E\approx3.3$--$3.9$.
\todo{Confirm the interpretation with Prof.\ Turaev.}

\subsection{The circular family}\label{sec:circular}
Newton started directly from Case~3 stalls, but it converges from a nearby
symmetric guess at the same energy, $\q_1=(3.6,3.6)$, using the diagonal
section $x_1=y_1$. Following this orbit upward in energy
(Figure~\ref{fig:circular}), Newton converges at all 21 energies from
$E\approx2.01$ to $4.02$, and every computed orbit is linearly stable
(Figure~\ref{fig:multipliers}(b)). Two of the three multiplier pairs coincide
at every energy, presumably because of the symmetry of the orbit. Near the
bottom of this range the remaining pair moves quickly around the unit circle
and comes close to $+1$ ($E\approx2.11$) and to $-1$ ($E\approx2.31$), so with
energy steps of $0.1$ there may be narrow instability intervals that we have
not resolved. \todo{Check with Prof.\ Turaev whether the double pair is
expected from the symmetry.} As the energy
increases, the loop grows towards the walls: its crossing of the diagonal moves
from $(3.59,3.59)$ to $(2.59,2.59)$. We then corrected the orbit at much higher
energies, $E\approx52$ and $E\approx202$ (with $h=10^{-5}$ in the latter case).
There it has become a diamond that touches the middle of each wall, and it is
still linearly stable. So, among the orbits we tested, this family is the most
robust, at least at the energies we sampled.

\begin{figure}[H]
  \centering
  \includegraphics[width=0.8\textwidth]{fig_circular_family.pdf}
  \caption{One period of the circular family; the two particles move on the
  same curve, half a period apart. (a) For $2.0\le E\le4.0$ the loop grows
  towards the walls. (b) At $E\approx52$ and $202$ it becomes a diamond
  touching all four walls. Every orbit shown is linearly stable.}
  \label{fig:circular}
\end{figure}

\subsection{Phase-space portrait at $E\approx52$}\label{sec:lyapunov}
Figure~\ref{fig:lyapunov} shows the Poincar\'e section $y_1=5$, $p_{y_1}>0$ on
$\M$ at $E\approx52.23$, the energy of the high-energy circular orbit, with
every point coloured by the finite-time Lyapunov exponent of its trajectory.
The phase space is clearly mixed. A large chaotic sea, where
$\lambda_{\max}\approx0.15$--$0.27$, surrounds many islands of regular motion
with $\lambda_{\max}$ close to zero. The histogram of $\lambda_{\max}$ is
correspondingly bimodal, with a sharp peak near $0$ and a broad one around
$0.2$. The largest islands are centred on the axis $p_{x_1}=0$ and come in
pairs, because of the reflection symmetries of the section.

Figure~\ref{fig:islands} shows the elliptic periodic orbits found at island
centres. Orbit~\#1, which crosses the section once per period ($k=1$,
$T\approx2.86$), is the vertical bouncing orbit of Section~\ref{sec:vertical},
now at much higher energy. Orbits \#2--\#10 cross the section two to four times
per period and have more complicated shapes, with periods up to $T\approx31$.
They are the kind of stable periodic motion that, by \cite{turaev2022}, keeps
the system from being ergodic at high energy. \todo{Check that this connection
is fair, and that \#1 really is the continuation of the vertical family.}

\begin{figure}[H]
  \centering
  \includegraphics[width=\textwidth]{lyapunov_poincare_section_reduced_report.jpg}
  \caption{Poincar\'e section $y_1=5$, $p_{y_1}>0$ on the centre-symmetric
  set $\M$ at $E\approx52.23$, with each crossing coloured by the finite-time
  maximal Lyapunov exponent of its trajectory ($T=500$, 13\,920 trajectories).
  Right: histogram of $\lambda_{\max}$. Dark points (small $\lambda_{\max}$)
  form islands of regular motion in the chaotic sea.}
  \label{fig:lyapunov}
\end{figure}

\begin{figure}[H]
  \centering
  \includegraphics[width=\textwidth]{lyapunov_section_periodic_orbits_report.jpg}
  \caption{Elliptic periodic orbits at island centres, $E\approx52.23$.
  Top: their points on the section. Below: the orbits in the box (particle~1
  blue, particle~2 orange; the black dot is the section point), with the number
  $k$ of section crossings per period and the period $T$.}
  \label{fig:islands}
\end{figure}

% ---------------------------------------------------------------------------
\section{Discussion and limitations}\label{sec:discussion}

The overall picture from these experiments is as follows. Near the equilibrium
the motion is regular. Slightly above it, $E\approx2.03$--$3$, chaos is
widespread. At higher energy, regular motion reappears around a small number of
linearly stable periodic orbits: the vertical bouncing orbit, the circular
orbit, and the island centres of Figure~\ref{fig:islands}. The vertical orbit
loses stability as the energy decreases, apparently through a symmetry-breaking
bifurcation, while the circular orbit is linearly stable at all the energies
we sampled, over two orders of magnitude.

\paragraph{Where continuation failed.}
Following the circular family \emph{downwards} from the Case~3 energy $E_0$
was much harder than following it upwards. With steps of $0.1\%$ in energy,
Newton still converged at $0.999E_0$, $0.998E_0$ and $0.997E_0$, but the
period of the computed orbit jumped from $101.8$ to $109.2$, $135.4$ and
$140.6$. In addition, at $0.998E_0$ the orbit was no longer symmetric and was
unstable (a multiplier $\lambda\approx-5.27$). So Newton was converging to
\emph{different} nearby periodic orbits rather than following one branch. At
$0.996E_0$ Newton failed: the residual stalled at about $0.06$. Retrying with
half the energy step, Newton converged both at the midpoint and then at
$0.996E_0$, so the immediate cause of the failure was the step size. Along the
converged sequence we also recorded the distance of the multipliers from $1$,
$\min_i|\lambda_i-1|$ (a fold would drive it to zero), and the condition number
of $DP-I$; they went from $1.56$ to $0.20$ and from $10$ to $216$ respectively.
This trend may indicate a multiplier approaching $1$, but because the orbits
along the sequence are not one branch, we cannot draw a conclusion from it.
Our interpretation is that several periodic orbits lie close together just
below $E_0$, and that natural-parameter continuation with this step size
cannot tell them apart. A method that
follows a branch along its arclength, such as pseudo-arclength continuation
\cite{kuznetsov}, together with a check that the symmetry is kept, would be the
natural next step.

\paragraph{Numerical reliability.}
The periodic orbits are well resolved: residuals are at most $10^{-9}$, and the
periods converge at second order under step halving. Multipliers are less
accurate (about $10^{-6}$), since the Jacobians use finite differences; this
does not affect the stable/unstable classification except very close to a
bifurcation. Finite-time Lyapunov exponents only estimate the true exponents.
Very thin chaotic layers, or very slow chaotic diffusion, may be invisible at
$T=500$.

\paragraph{Scope.}
The orbits of Sections~\ref{sec:vertical}--\ref{sec:circular} lie on $\M$,
but their multipliers are computed on the full six-dimensional section, so
their linear stability includes perturbations that leave $\M$. In contrast,
the Lyapunov exponents and island centres of Section~\ref{sec:lyapunov} are
computed within $\M$ only; an orbit that is elliptic there could still be
unstable in directions transverse to $\M$. \todo{re-run the transverse check in
\texttt{lyapunov\_section\_periodic\_orbits.py} and report the result.}

\paragraph{Open questions.}
\begin{itemize}
  \item Resolve the bifurcation of the vertical orbit near $E\approx3.3$--$3.9$
        and follow the S-shaped orbits that appear.
  \item How does the size of the chaotic sea change with energy at high energy,
        and does the stable region grow as \cite{turaev2022} suggests?
  \item How do the results depend on the exponents $\alpha,\beta$ and on the
        shape of the box?
\end{itemize}

% ---------------------------------------------------------------------------
\section*{Code and reproducibility}
All code is available at
\url{https://github.com/jiajunwu1729/2particle-numerics}
\todo{(tag/commit used for this report)}. Table~\ref{tab:repro} lists the
script that produces each figure.

\begin{table}[H]
\centering
\caption{Figures and the code that produces them.}
\label{tab:repro}
\begin{tabular}{@{}ll@{}}
\toprule
Figure & Source\\
\midrule
\ref{fig:cases}, \ref{fig:energy-sections} & \texttt{report/make\_section\_figures.py} (same computation as \texttt{report03.ipynb})\\
\ref{fig:vertical}, \ref{fig:circular} & \texttt{report/make\_orbit\_figures.py} (periodic points from \texttt{report04.ipynb})\\
\ref{fig:multipliers} & \texttt{report/make\_multiplier\_figure.py} (multipliers from \texttt{report04.ipynb})\\
\ref{fig:lyapunov} & \texttt{lyapunov\_poincare\_section.py}\\
\ref{fig:islands} & \texttt{lyapunov\_section\_periodic\_orbits.py}\\
& (both reduced for the PDF by \texttt{report/compress\_figures.py})\\
\bottomrule
\end{tabular}
\end{table}

\section*{Acknowledgements}
I thank Prof.\ Dmitry Turaev for suggesting this project and for his guidance.
\todo{funding, if any.}

% ---------------------------------------------------------------------------
\begin{thebibliography}{9}
\bibitem{turaev2022}
V.~Rom-Kedar and D.~Turaev,
Stable motions of high energy particles interacting via a repelling potential,
arXiv:2208.14993v2, 2023. \todo{check whether it has since been published}

\bibitem{hairer2006}
E.~Hairer, C.~Lubich and G.~Wanner,
\emph{Geometric Numerical Integration}, 2nd ed., Springer, 2006.

\bibitem{kuznetsov}
Yu.~A.~Kuznetsov, \emph{Elements of Applied Bifurcation Theory}, 3rd ed.,
Springer, 2004.

\bibitem{meiss}
J.~D.~Meiss, \emph{Differential Dynamical Systems}, revised ed., SIAM, 2017.
\end{thebibliography}

\end{document}
